\section{Setting}
\label{sec:setting}

\subsection{Yukon}

Yukon describes itself as a platform for open frontier research, on which organisers create autoresearch challenges with verifiable benchmarks and solvers deploy agents against them \citep{yukon2026}.
Each benchmark has a public repository, a set of editable paths, a scoring direction, a baseline score, and an evaluator that scores each submission.
The public API at \url{https://api.yukon.org} lists every benchmark, open or closed, and returns each benchmark's submissions without credentials.

A solver submits through the Yukon command-line tool.
The tool requires a public Markdown note, a \field{-{}-model} flag naming the underlying model, and a \field{-{}-harness} flag naming the coding agent or harness \citep{yukoncli2026}.
It writes these two flags into the note as canonical \field{Model:} and \field{Harness:} lines.
The tool's guidance asks for the exact, fully qualified model and tells solvers never to copy the model attribution from the base submission \citep{yukoncli2026}.
It also asks solvers to cite a promoted submission they started from in the note, and to credit substantial unpromoted help from another solver through a co-author flag.

The evaluator either accepts or rejects a submission, or the run fails.
An accepted submission that improves on the current best is promoted onto the benchmark branch, and the landed commit is recorded as \field{promotedSourceRef} \citep{yukoncli2026}.
Each benchmark object carries a field \field{minScoreImprovementBips}, the smallest relative improvement that can promote, and the field differs between benchmarks.
Normal synchronisation restores the editable paths from the best promoted submission, so the current best is the default starting point for new work \citep{yukoncli2026}.
A separate endpoint, \field{/api/submissions/\{id\}/attribution}, returns a structured attribution record for a submission.
When the solver declared parents, the record has \field{declaration} set to \field{provided} and lists them as \field{dependencies}.
Both fields are described here from the API responses, as no public documentation of them was found.

\subsection{ECDSA.Fail}

ECDSA.Fail, launched by Eigen Labs in late May 2026, optimises a reversible secp256k1 mixed point-addition circuit and scores it by $S = Q \times T$, the product of peak logical qubit width and average executed Toffoli count \citep[\S\S3, 4.2]{long2026ecdsafail}.
Over about eight weeks, more than one hundred participants using diverse models produced over four hundred promoted submissions \citep[\S4.2]{long2026ecdsafail}.
Promoted improvements add their score, source, participant attribution and declared AI model to the public record \citep[\S4.2]{long2026ecdsafail}.
On Yukon the benchmark is \bench{gpsanant/ecdsafail-challenge}.
It is the benchmark that motivated this study, so it is always included and its own primary estimate is reported as the headline case.

\subsection{Study benchmarks}

The pre-registration fixes a selection rule, applied by \file{analysis/study.py} to the frozen snapshot and the frozen lineage files.
A benchmark joins the study, open or closed, when it meets all of the following.
\begin{enumerate}
  \item It has at least 500 analysis rows (Section~\ref{sec:denominator}), and at least 100 of them were promoted.
  \item At least two model families each hold at least 10\% of the analysis rows that carry a single known family.
  \item Its trees are recoverable: a confirmatory window exists (Section~\ref{sec:window}), and at least half of its promoted analysis rows have a tree commit that the lineage build did not record as unfetchable.
\end{enumerate}
ECDSA.Fail is included whatever the rule says.

Ten benchmarks met this rule at a probe on 17 October 2026, and attribution records were frozen for all ten.
Table~\ref{tab:benchmarks} lists them with their row counts in the frozen snapshot.
The rule ran again on the frozen snapshot and lineage before the confirmatory analysis and returned the same \StudyK{} benchmarks.
Two pairs are siblings: the two flock benchmarks, and the two EIP-8200 benchmarks, which share one repository.
Sensitivity analysis 8 reduces each pair to one member (Section~\ref{sec:sensitivity}).

\begin{table}[t]
\centering
\caption{Study benchmarks. Rows are all submissions in the frozen snapshot of 18 October, cancelled included; attribution 404 counts the rows whose attribution endpoint returned HTTP 404 at the freeze. Source: \file{snapshot/FREEZE.md}.}
\label{tab:benchmarks}
\begin{tabular}{lrr}
\toprule
Benchmark & Rows & Attribution 404 \\
\midrule
\bench{davidtai/mlxfast-gemma4-26b-a4b} & 3,239 & 397 \\
\bench{davidtai/qwen38-125b-a6b-cuda-v1} & 5,757 & 337 \\
\bench{eigenlabs/eip8200-challenges/modexp} & 1,512 & 151 \\
\bench{eigenlabs/eip8200-challenges/ripemd160} & 1,233 & 122 \\
\bench{eigenlabs/flock-challenge} & 4,632 & 53 \\
\bench{eigenlabs/flock-challenge-multi/x86} & 3,392 & 166 \\
\bench{eigenlabs/lighter-prover-challenge} & 8,379 & 338 \\
\bench{eigenlabs/mlxfast-challenge} & 1,904 & 85 \\
\bench{eigenlabs/sig-golf/full} & 933 & 246 \\
\bench{gpsanant/ecdsafail-challenge} & 1,593 & 255 \\
\midrule
Total & 32,574 & 2,150 \\
\bottomrule
\end{tabular}
\end{table}
